% 2.7 — Synthèse et positionnement
\section{Synthèse et positionnement de \texorpdfstring{\demo}{demo\_cli}}
\label{sec:synthese}

Le tableau~\ref{tab:existant} compare les approches selon six critères tirés de la
problématique.

\newcommand{\oui}{\textcolor{dALLOW}{\textbf{oui}}}
\newcommand{\non}{\textcolor{dDENY}{non}}
\newcommand{\partl}{\textcolor{dASK}{partiel}}
\begin{table}[htbp]
\centering
\caption{Les protections existantes face aux critères du problème}
\label{tab:existant}
\small
\begin{tabularx}{\linewidth}{@{}L{3.5cm}YYYYYY@{}}
\toprule
& \textbf{Prédit l'effet} & \textbf{Récupère} & \textbf{Couvre le shell} &
  \textbf{Traite le distant} & \textbf{Sans friction} & \textbf{Linux et Windows} \\
\midrule
Demande de permission & \non & \non & \oui & \partl{} (la personne) & \non & \oui \\
Juge automatique & \partl{} & \non & \oui & \partl{} & \oui & \oui \\
Bac à sable & \non & \non & \oui & \partl{} (réseau) & \oui & \partl{} \\
Points de restauration de l'agent & \non & \partl{} (édition) & \non & \non & \oui & \oui \\
Git, sauvegardes & \non & \partl{} & \partl{} & \non & \oui & \oui \\
Motifs (regex) & \partl{} & \non & \oui & \partl{} & \oui & \oui \\
Filet d'exécution seul & \non & \oui{} (local) & \oui & \non & \oui & \partl{} \\
\midrule
\textbf{\demo{}} & \oui & \oui & \oui & \oui{} (ask/deny) & \partl{} & \oui \\
\bottomrule
\end{tabularx}
\end{table}

Aucune protection existante ne réunit les deux propriétés centrales : \emph{prédire}
l'effet d'une ligne et \emph{récupérer} ce qu'elle détruit. \demo{} se positionne donc
comme une couche qui \emph{s'ajoute} aux autres, avec quatre partis pris :
\begin{enumerate}
  \item \textbf{Prédire par la structure et la simulation}, pas par les motifs.
  \item \textbf{Récupérer d'abord}, ne demander que ce qui ne peut pas être récupéré, et
    refuser ce dont il n'existe aucun retour réaliste.
  \item \textbf{Un filet d'exécution comme plancher} pour ce qui ne peut pas être lu.
  \item \textbf{Ne jamais revendiquer une protection non prise} : chaque décision dit ce
    qui a été capturé, et ce qui ne l'a pas été.
\end{enumerate}
La case « sans friction » reste partielle : une escalade interrompt l'agent, et c'est
voulu pour ce qui ne peut pas être récupéré.

\section*{Conclusion du chapitre}
Les protections existantes sont réelles mais orthogonales au problème : les demandes de
permission s'usent, les bacs à sable laissent le projet inscriptible, les points de
restauration ignorent le shell, et l'analyse syntaxique s'arrête à la structure. Le
chapitre suivant décrit l'architecture qui comble ce vide : comment \demo{} lit, simule,
décide, récupère et surveille.
